\documentclass[10pt,a4paper]{amsart}
\usepackage{amsmath,amssymb,amsthm,graphicx}
% Retained preamble for the Stacks Project chapter {CH} (shared project preamble).
% The chapter file pinned-original-source.tex only carries \input{{preamble}} on its first line,
% so no preamble text lives inside the chapter itself.  The Stacks Project uses one single shared
% preamble for every chapter: the snapshot stages it three times (sources/stacks-chow/preamble.tex,
% sources/stacks-etale-cohomology/preamble.tex and focused-seeds/AGP/sources/stacks-snapshot-20261011/
% preamble.tex) and those copies are identical in size.  The block below is copied mechanically,
% line for line, from sources/stacks-chow/preamble.tex lines 14-18 and 154-217: the Xy-pic setup and
% the native theorem declarations together with the mathematical macros.  Nothing was added, removed
% or redefined; the wrapper class and the amsmath/amssymb/amsthm/graphicx base packages are supplied
% by the assembler.
% For commutative diagrams we use Xy-pic
% Chapter: stacks-cohomology.tex -- Cohomology of Sheaves.
\usepackage[all]{xy}

% We use 2cell for 2-commutative diagrams.
\xyoption{2cell}

% Theorem environments.
%
\theoremstyle{plain}
\newtheorem{theorem}[subsection]{Theorem}
\newtheorem{proposition}[subsection]{Proposition}
\newtheorem{lemma}[subsection]{Lemma}

\theoremstyle{definition}
\newtheorem{definition}[subsection]{Definition}
\newtheorem{example}[subsection]{Example}
\newtheorem{exercise}[subsection]{Exercise}
\newtheorem{situation}[subsection]{Situation}

\theoremstyle{remark}
\newtheorem{remark}[subsection]{Remark}
\newtheorem{remarks}[subsection]{Remarks}

\numberwithin{equation}{subsection}

% Macros
%
\def\lim{\mathop{\mathrm{lim}}\nolimits}
\def\colim{\mathop{\mathrm{colim}}\nolimits}
\def\Spec{\mathop{\mathrm{Spec}}}
\def\Hom{\mathop{\mathrm{Hom}}\nolimits}
\def\Ext{\mathop{\mathrm{Ext}}\nolimits}
\def\SheafHom{\mathop{\mathcal{H}\!\mathit{om}}\nolimits}
\def\SheafExt{\mathop{\mathcal{E}\!\mathit{xt}}\nolimits}
\def\Sch{\mathit{Sch}}
\def\Mor{\mathop{\mathrm{Mor}}\nolimits}
\def\Ob{\mathop{\mathrm{Ob}}\nolimits}
\def\Sh{\mathop{\mathit{Sh}}\nolimits}
\def\NL{\mathop{N\!L}\nolimits}
\def\CH{\mathop{\mathrm{CH}}\nolimits}
\def\proetale{{pro\text{-}\acute{e}tale}}
\def\etale{{\acute{e}tale}}
\def\QCoh{\mathit{QCoh}}
\def\Ker{\mathop{\mathrm{Ker}}}
\def\Im{\mathop{\mathrm{Im}}}
\def\Coker{\mathop{\mathrm{Coker}}}
\def\Coim{\mathop{\mathrm{Coim}}}

% Boxtimes
%
\DeclareMathSymbol{\boxtimes}{\mathbin}{AMSa}{"02}

%
% Macros for moduli stacks/spaces
%
\def\QCohstack{\mathcal{QC}\!\mathit{oh}}
\def\Cohstack{\mathcal{C}\!\mathit{oh}}
\def\Spacesstack{\mathcal{S}\!\mathit{paces}}
\def\Quotfunctor{\mathrm{Quot}}
\def\Hilbfunctor{\mathrm{Hilb}}
\def\Curvesstack{\mathcal{C}\!\mathit{urves}}
\def\Polarizedstack{\mathcal{P}\!\mathit{olarized}}
\def\Complexesstack{\mathcal{C}\!\mathit{omplexes}}
% \Pic is the operator that assigns to X its picard group, usage \Pic(X)
% \Picardstack_{X/B} denotes the Picard stack of X over B
% \Picardfunctor_{X/B} denotes the Picard functor of X over B
\def\Pic{\mathop{\mathrm{Pic}}\nolimits}
\def\Picardstack{\mathcal{P}\!\mathit{ic}}
\def\Picardfunctor{\mathrm{Pic}}
\def\Deformationcategory{\mathcal{D}\!\mathit{ef}}
\title{The Cech complex of pushforwards of the restrictions is a resolution of the O\_X-module}
\author{The Stacks Project Authors}
\date{stacks-project:cohomology.tex:local-snapshot-20261011; GNU FDL 1.2 or later}
\begin{document}
\maketitle
\noindent\textit{Editorial note.} Preamble: only the retained preamble of this chapter (sources/stacks-cohomology/retained-preamble.tex) is taken, because the chapter file itself carries just an input of the project's single shared preamble on its first line. It holds the Xy-pic setup, the author's native theorem declarations and the mathematical macros, copied line for line; nothing was added, removed or redefined. This card uses neither stix nor transparent, which the host TeX Live lacks.\par\medskip
\noindent\textit{Editorial note.} The display of the Cech differential (equation-d-cech, source lines 895-908) is printed as a context span, so that the differential formula which the counted statement names is visible in the material. Consequently this card needs no external newlabel stub at all; both local trial XeLaTeX runs log no undefined reference.\par\medskip
\noindent\textit{Editorial note.} Original numbering: in this chapter lemmas, equations and remarks share the subsection counter, and material-only assembly locates passages by label, so the printed numbers need not agree with the original ones.\par\medskip
\noindent\textit{Editorial note.} Uniform work/source\_title across every card of this source, as section 11 of NEW\_TASK\_AUTHORING.md requires (the manifest field work is taken from source\_title). The precise section of each claim is carried by claim\_key and by statement\_anchor in the registration, and the section name is repeated in the editorial notes where it matters.\par\medskip
\noindent\textit{Allowed standard prerequisites (editorial).}\begin{itemize}
\item The Cech complex of an open covering of a sheaf of modules and its coboundary formula. Same chapter, the display labelled equation-d-cech (canonical Stacks tag 01EE, source line 902, shown as context above); the counted statement names it for the differential.. Fixes the terms C\textasciicircum{}p(U,F) = prod\_\{i\_0...i\_p\} (j\_\{i\_0...i\_p\})\_* F\_\{i\_0...i\_p\} and the differential d whose exactness on stalks the proof verifies.
\item Exactness of a sequence of sheaves of modules can be checked on stalks, and restriction to an open subspace commutes with taking stalks. Standard sheaf theory as used throughout the same chapter's section The Cech complex and Cech cohomology (source lines 868 onwards); the proof uses it at lines 4603-4609.. Reduces the quasi-isomorphism claim to exactness of the extended complex of stalks, which is then established by an explicit contracting homotopy.
\end{itemize}\medskip
$$
d : \check{\mathcal{C}}^p(\mathcal{U}, \mathcal{F})
\longrightarrow
\check{\mathcal{C}}^{p + 1}(\mathcal{U}, \mathcal{F})
$$
by the formula
\begin{equation}
\label{equation-d-cech}
d(s)_{i_0\ldots i_{p + 1}}
=
\sum\nolimits_{j = 0}^{p + 1}
(-1)^j
s_{i_0\ldots \hat i_j \ldots i_{p + 1}}|_{U_{i_0\ldots i_{p + 1}}}
\end{equation}

\begin{lemma}
\label{lemma-covering-resolution}
Let $X$ be a ringed space. Let $\mathcal{U} : X = \bigcup_{i \in I} U_i$
be an open covering of $X$. Let $\mathcal{F}$ be an $\mathcal{O}_X$-module.
Denote $\mathcal{F}_{i_0 \ldots i_p}$ the restriction of
$\mathcal{F}$ to $U_{i_0 \ldots i_p}$. There exists a complex
${\mathfrak C}^\bullet(\mathcal{U}, \mathcal{F})$
of $\mathcal{O}_X$-modules with
$$
{\mathfrak C}^p(\mathcal{U}, \mathcal{F}) =
\prod\nolimits_{i_0 \ldots i_p}
(j_{i_0 \ldots i_p})_* \mathcal{F}_{i_0 \ldots i_p}
$$
and differential
$d : {\mathfrak C}^p(\mathcal{U}, \mathcal{F})
\to {\mathfrak C}^{p + 1}(\mathcal{U}, \mathcal{F})$
as in Equation (\ref{equation-d-cech}). Moreover, there exists a canonical
map
$$
\mathcal{F} \to {\mathfrak C}^\bullet(\mathcal{U}, \mathcal{F})
$$
which is a quasi-isomorphism, i.e.,
${\mathfrak C}^\bullet(\mathcal{U}, \mathcal{F})$
is a resolution of $\mathcal{F}$.
\end{lemma}

\begin{proof}
We check
$$
0 \to \mathcal{F} \to \mathfrak{C}^0(\mathcal{U}, \mathcal{F}) \to
\mathfrak{C}^1(\mathcal{U}, \mathcal{F}) \to  \ldots
$$
is exact on stalks. Let $x \in X$ and choose $i_{\text{fix}} \in I$
such that $x \in U_{i_{\text{fix}}}$. Then define 
$$
h : \mathfrak{C}^p(\mathcal{U}, \mathcal{F})_x
\to \mathfrak{C}^{p - 1}(\mathcal{U}, \mathcal{F})_x
$$
as follows: If $s \in \mathfrak{C}^p(\mathcal{U}, \mathcal{F})_x$, take
a representative
$$
\widetilde{s} \in
\mathfrak{C}^p(\mathcal{U}, \mathcal{F})(V) =
\prod\nolimits_{i_0 \ldots i_p}
\mathcal{F}(V \cap U_{i_0} \cap \ldots \cap U_{i_p})
$$
defined on some neighborhood $V$ of $x$, and set
$$
h(s)_{i_0 \ldots i_{p - 1}} =
\widetilde{s}_{i_{\text{fix}} i_0 \ldots i_{p - 1}, x}.
$$
By the same formula (for $p = 0$) we get a map
$\mathfrak{C}^{0}(\mathcal{U},\mathcal{F})_x \to \mathcal{F}_x$.
We compute formally as follows:
\begin{align*}
(dh + hd)(s)_{i_0 \ldots i_p}
& =
\sum\nolimits_{j = 0}^p
(-1)^j
h(s)_{i_0 \ldots \hat i_j \ldots i_p}
+
d(s)_{i_{\text{fix}} i_0 \ldots i_p}\\
& =
\sum\nolimits_{j = 0}^p
(-1)^j
s_{i_{\text{fix}} i_0 \ldots \hat i_j \ldots i_p}
+
s_{i_0 \ldots i_p}
+
\sum\nolimits_{j = 0}^p
(-1)^{j + 1}
s_{i_{\text{fix}} i_0 \ldots \hat i_j \ldots i_p} \\
& =
s_{i_0 \ldots i_p}
\end{align*}
This shows $h$ is a homotopy from the identity map of
the extended complex
$$
0 \to \mathcal{F}_x \to \mathfrak{C}^0(\mathcal{U}, \mathcal{F})_x
\to \mathfrak{C}^1(\mathcal{U}, \mathcal{F})_x \to \ldots
$$
to zero and we conclude.
\end{proof}

\end{document}
